\chapter{Bloc I : Anàlisi}
Tot el temari d'aquest bloc ja s'ha treballat a Matemàtiques II (MAT). Com que aquest manual està pensat per estudiants que han cursat MAT, l'alumne hauria de ser capaç de resoldre aquests problemes sense dificultats, així que no s'explicaran els temes del bloc I.
\section{Àlgebra lineal}
D'àlgebra lineal tal com està en el document del currículum de l'examen de les Matemàtiques Aplicades a les Ciències Socials \footnote{\url{https://universitats.gencat.cat/web/.content/01_acces_i_admissio/proves-acces-PAU-PAP/ambits/matematiques-cs/02_Pau25_MACS_orientacions.pdf}} hi ha:
\begin{itemize}
 \item Operacions amb matrius: suma, producte. Equacions matricials ($2 \times 2$).
 \item Resolució de sistemes d'equacions sense paràmetres. Problemes que es poden modelitzar mitjançant sistemes d'equacions.
\end{itemize}

\section{Càlcul}
Segons indica el mateix currículum de l'examen de les MACS entren els temes següents de càlcul:
\begin{itemize}
 \item Identificació, representació i anàlisi de les propietats de funcions polinòmiques, exponencials, logarítmiques i funcions a trossos.
 \item Càlcul de límits de funcions racionals. Aplicació al càlcul d'asímptotes horitzontals i verticals d'aquestes funcions.
 \item Càlcul de derivades de funcions. Estudi local d'una funció: continuïtat, derivabilitat, extrems relatius. Creixement i decreixement d'una funció.
 \item Resolució de problemes d'optimització.
\end{itemize}
